\documentclass[11pt]{article}
\usepackage[letterpaper,left=.9in,right=.9in,top=.8in,bottom=.82in]{geometry}
\usepackage{amsmath,amssymb,amsthm,xcolor,hyperref}
\hypersetup{colorlinks=true,linkcolor=blue!45!black,urlcolor=blue!45!black,citecolor=blue!45!black,
pdftitle={A third-order bound for unions of odd residue classes},pdfauthor={GitHoobar}}
\newtheorem{theorem}{Theorem}
\newtheorem{lemma}[theorem]{Lemma}
\newcommand{\D}{\mathcal D}
\title{\vspace{-1.3em}A third-order bound for unions of odd residue classes}
\author{GitHoobar}
\date{27 November 2026}
\begin{document}
\maketitle
\begin{abstract}
For every nontrivial divisor of an odd integer \(L\), choose one residue
class. The size of their union is at most the sum of their sizes, less the
pairwise-coprime intersection terms, plus the pairwise-coprime triple
intersection terms. The bound holds for every number of prime factors
and every choice of their exponents. This answers the unrestricted
third-order question distinguished by Toker from the all-orders
CRT coverage benchmark. The key step is to show that whenever the
third-order bound is below \(L\), all induced independence polynomials
are positive and their full alternating sums dominate their cubic truncations
on complete divisor families.
\end{abstract}

\section{The coverage bound}
Let \(L\) be an odd positive integer and let \(\D=\{d>1:d\mid L\}\).
Choose a residue \(a_d\bmod d\) for each \(d\in\D\). Write \(U\) for the
number of residues modulo \(L\) lying in their union. For \(j\geq 1\) put
\[
 A_j=\sum_{\substack{J\subseteq\D,\ |J|=j\\
               \gcd(d,e)=1\ \text{for distinct }d,e\in J}}
               \ \prod_{d\in J}\frac1d,\qquad A_0=1 .
\]
All families in this definition are unordered. By the Chinese remainder
theorem (CRT), each summand is the fraction of the period lying in the
specified intersection, for every choice of residues.

\begin{theorem}\label{thm:main}
For every \(L\) and every assignment as above,
\[
 \frac{U}{L}\le b,\qquad b=A_1-A_2+A_3 .
\]
If \(b<1\), then also
\[
 1-\frac UL\ \geq\ \sum_{j\ge0}(-1)^j A_j\ \geq\ 1-b>0.
\]
\end{theorem}

Thus at least \(L(1-b)\) points are uncovered whenever \(b<1\).
The usual third-order Bonferroni upper bound includes all actual pair
and triple intersections. Here only pairwise-coprime families enter,
so the stated upper bound depends on \(L\) alone.
Hackmon introduced the third-order expression \cite[Definition 2]{hackmon}.
Toker distinguished it from the all-orders expression, established the
third-order inequality through seven prime factors, and disproved the
unrestricted all-orders bound \cite[\S\S 4--5]{toker}.
Theorem~\ref{thm:main} settles the third-order question at unrestricted rank.
It gives a criterion for individual periods. The existence of a finite
cover with distinct odd moduli remains open.

\section{Two inequalities on a cube}
Factor \(L=\prod_{i=1}^k p_i^{a_i}\), and define
\[
 x_i=\sum_{e=1}^{a_i}p_i^{-e},\qquad P(t)=\prod_{i=1}^k(1+t x_i),
 \qquad W=A_1=P(1)-1 .
\]
Oddness gives \(0\leq x_i\leq 1/(p_i-1)\leq 1/2\).
In an ordered list of pairwise-coprime divisors, a prime coordinate
can go to at most one slot. For \(t\) slots, each allowed exponent for
that prime contributes weight \(p_i^{-e}\), so the product \(P(t)\)
counts all weighted placements with empty slots allowed.
Excluding empty slots and then dividing by \(j!\) gives
\[
 j! A_j=\sum_{t=0}^{j}(-1)^{j-t}\binom jt P(t).
\]
In particular \cite[eq.~(5.2)]{toker},
\begin{equation}\label{eq:b}
 6b=P(3)-6P(2)+15P(1)-10 .
\end{equation}

\begin{lemma}\label{lem:cube}
On every cube \([0,1/2]^k\), the function \(b\) is nondecreasing in each
coordinate and \(W\leq 5b\).
\end{lemma}
\begin{proof}
A function affine separately in each argument is a convex combination
of its values at the vertices of the cube, by successive linear
interpolation. The slope of \(6b\) in a given coordinate, writing \(P_*\)
for the product over the remaining coordinates, is
\(3P_*(3)-12P_*(2)+15P_*(1)\).
At a vertex with \(m\) coordinates equal to \(1/2\), its sign is that of
\[
 a_m=3\cdot 5^m-12\cdot4^m+15\cdot3^m.
\]
For \(m\ge7\) use \((5/4)^7>4\). The smaller cases appear below.
Likewise \(6(5b-W)=5P(3)-30P(2)+69P(1)-44\) is separately affine.
Its vertex values have the signs of
\[
 c_m=5\cdot5^m-30\cdot4^m+69\cdot3^m-44\cdot2^m .
\]
For \(m\ge9\) use \((5/4)^9>6\) and \(69\cdot3^m>44\cdot2^m\).
The remaining integer checks are
\[
\begin{array}{c|rrrrrrrrr}
m   &0&1&2&3&4&5&6&7&8\\\hline
a_m &6&12&18&12&18&732&8658&\cdot&\cdot\\
c_m &0&24&90&216&330&264&2730&44376&428490 .
\end{array}
\]
These two vertex inequalities prove both assertions.
\end{proof}

\section{From the cubic truncation to a probability bound}
For \(S\subseteq\D\), let \(A_j(S)\) count the same weighted families
using only divisors in \(S\), and put \(Q(S)=\sum_{j\ge0}(-1)^j A_j(S)\).
The sum is finite. Connect two divisors by an edge when they share a
prime. Write \(N[v]\) for \(v\) and all its neighbors in this graph.

\begin{lemma}\label{lem:pos}
If \(b<1\), then \(Q(\D)\geq1-b\) and \(Q(S)>0\) for every \(S\subseteq\D\).
\end{lemma}
\begin{proof}
For a set \(I\) of prime coordinates, let \(\D_I\) contain all
nontrivial divisors supported in \(I\), with the original exponent
bounds. Its cubic bound \(b_I\) is obtained by zeroing the coordinates
outside \(I\). Lemma~\ref{lem:cube} gives \(b_I\leq b<1\) and
\(W_I=A_1(\D_I)\leq5b_I<5\).
Writing \(A_j^I=A_j(\D_I)\), count an admissible family of size \(j+1\)
together with a marked member. The member that is added to a family of
size \(j\) has total available weight at most \(W_I\), so
\[
 (j+1)A_{j+1}^I\leq W_I A_j^I .
\]
Thus \(A_{j+1}^I\leq A_j^I\) for \(j\ge4\). Pairing the tail
\((A_4^I-A_5^I)+(A_6^I-A_7^I)+\cdots\) gives
\(Q(\D_I)\geq 1-b_I>0\). This includes \(I\) equal to all coordinates.

To pass to arbitrary \(S\subseteq\D_I\), induct on \(|I|\). The empty
coordinate set gives \(Q(\varnothing)=1\). Start at \(\D_I\) and delete
unwanted vertices one at a time. Splitting independent sets according to
whether they contain \(v\in T\) gives
\begin{equation}\label{eq:delete}
 Q(T\setminus\{v\})=Q(T)+\frac1v Q(T\setminus N[v]).
\end{equation}
Every element of \(T\setminus N[v]\) avoids every prime of \(v>1\).
The induction hypothesis therefore gives \(Q(T\setminus N[v])>0\),
regardless of earlier deletions. Each deletion preserves positivity.
\end{proof}

\begin{proof}[Proof of Theorem~\ref{thm:main}]
The conclusion for \(b\ge1\) follows from \(U\le L\). Suppose \(b<1\).
For \(T\subseteq\D\), let \(Z(T)\) be the probability that a uniform point
modulo \(L\) avoids the classes in \(T\). We include the short ratio
argument for Shearer's inequality \cite{shearer}. The class at \(v\)
is independent of the joint avoidance event for \(R=T\setminus N[v]\)
by CRT. With \(T'=T\setminus\{v\}\), this gives
\[
 Z(T)\ \geq\ Z(T')-\frac1v Z(R).
\]
We prove inductively on \(|T|\) that \(Z(T)>0\) and
\(Z(T)/Z(T')\geq Q(T)/Q(T')\) for each \(v\in T\).
Both functions equal \(1\) at \(\varnothing\). The inductive ratio bounds
telescope along a chain from \(R\) to \(T'\), giving
\(Z(T')/Z(R)\geq Q(T')/Q(R)\), including equality on an empty chain.
All denominators are positive by induction and Lemma~\ref{lem:pos}.
Consequently
\[
 \frac{Z(T)}{Z(T')}\geq 1-\frac1v\,\frac{Q(R)}{Q(T')}
   =\frac{Q(T)}{Q(T')}>0 ,
\]
by \eqref{eq:delete}. This proves the induction. Telescoping again gives
\(1-U/L=Z(\D)\geq Q(\D)\geq1-b\), as required.
\end{proof}

\paragraph{A range where the bound gives information.}
For \(L=3\cdot5\cdot7\cdot11\cdot13=15015\), the single-class sum is
\(LW=17241>L\), while \(Lb=12110<L\).
For any prescribed number of additional prime factors, adjoining
sufficiently large distinct primes preserves both strict inequalities by
continuity of the products \(P(t)\) at zero. Thus the estimate can certify
uncovered points beyond the raw sum of class sizes at arbitrarily large rank.

\paragraph{Provenance.}
This research note was developed with OpenAI Codex under the GitHub
handle GitHoobar. The benchmark and its distinction from the all-orders
variant are due to the sources cited below. The argument is presented
for independent review.

\begin{thebibliography}{9}
\bibitem{hackmon} I.~Hackmon, \emph{On odd covering systems with few prime
factors}, preprint, April 2026. \href{https://rxiv.org/abs/2604.0035}{viXra:2604.0035}.
\bibitem{toker} S.~A.~Toker, \emph{On the CRT Coverage Maximality Conjecture
and the Erd\H{o}s--Selfridge Problem Through Rank Seven}, research memorandum,
3 November 2026. \href{https://doi.org/10.5281/zenodo.23927150}{doi:10.5281/zenodo.23927150}.
\bibitem{shearer} J.~B.~Shearer, \emph{On a problem of Spencer},
Combinatorica \textbf{5} (1985), 241--245.
\href{https://doi.org/10.1007/BF02579368}{doi:10.1007/BF02579368}.
\end{thebibliography}
\end{document}
